\section{结论}
\label{sec:conc}

本文重新审视了电力系统的分散式小增益--相位条件，并指出其用于 GFM 变流器时的固有局限。为克服这些限制，本文引入一种结合直角坐标导纳系和极坐标导纳系的一般化环路整形变换。该方法不仅解决了变流器的低频非扇形性，还能在整个频率范围内给出更紧的相位行为界。

本文首先在接入无限大母线的 GFM 变流器上演示了该方法，其中虚拟导纳滤波进一步降低了保守性；随后将方法应用于 IEEE 14 节点系统，并成功识别出导致失稳的设备。

尽管结果具有良好前景，仍有若干问题有待解决，特别是系统不仅包含 GFM 变流器，还包含同步机和 GFL 变流器等其他设备时。同步机与本文所研究的 GFM 变流器具有相似动态，因此预计所提方法仍然适用；但励磁器、自动电压调节器和电力系统稳定器等辅助控制环节的影响还需要进一步验证。

对于具有不同动态特性的 GFL 变流器，小增益条件更为关键，所提变换能够带来的收益也较有限。为了给 GFL 装置获得保守性更低的增益条件，一种有益做法是将 GFM 的高增益转移到网络侧，从而提高网络的等效强度。

未来工作将研究变换矩阵的系统化选择方法，并通过把网络拓扑信息嵌入变换来进一步降低保守性。
